% pspace.tex  — §6 of "Cocycle-Graded Separation Logic"
% Self-contained section; \input this file from the volume main.
% Requires: amsthm, amsmath, amssymb, hyperref (via the preamble).

\section{PSPACE-Completeness of Cocycle-Bounded Equivalence}
\label{sec:pspace}

We prove that $k$-cocycle program equivalence is
\textsc{PSPACE}-complete when $k$ is given in \emph{unary}.
Throughout, fix an idea monoid $\mathcal{M}=(I,{\circ},e,\mathit{rs})$
where $|I|=m$, $|\mathit{rs}|\le B$, and $\mathit{rs}$ has
\emph{pairing rank} $r$ (see Definition~\ref{def:rank}).
Program size is $n$.

\begin{theorem}[PSPACE-completeness, unary $k$]
\label{thm:pspace-main}
The problem
\[
  \textsc{CocycleEq} \;=\;
  \bigl\{\langle P,\,Q,\,\mathcal{M},\,k\rangle
        \mid P \sim_k Q\bigr\}
\]
is \textsc{PSPACE}-complete when $k$ is encoded in unary,
$\mathcal{M}$ is given as explicit multiplication and pairing tables,
and $P,Q$ are \emph{EffCalc} programs.
\end{theorem}

\noindent
The proof splits into a membership argument
(\S\ref{sec:pspace-upper}) and a hardness argument
(\S\ref{sec:pspace-lower}).

%-----------------------------------------------------------------
\subsection{The Quotient Monoid \texorpdfstring{$M_k$}{Mk}}
\label{sec:quotient}
%-----------------------------------------------------------------

\begin{definition}[Pairing rank]
\label{def:rank}
For $\mathcal{M}=(I,{\circ},e,\mathit{rs})$, the \emph{pairing rank}
$r$ is the rank of the integer matrix
$R \in \mathbb{Z}^{m\times m}$ with $R_{ad}=\mathit{rs}(a,d)$.
\end{definition}

\begin{definition}[$k$-indistinguishability]
\label{def:kindist}
$a \equiv_k b$ in $I$ iff for all $a',d\in I$:
\begin{align}
  |\mathit{rs}(a{\circ}a',d) - \mathit{rs}(b{\circ}a',d)| &\le k,
  \label{eq:c1}\\
  |\mathit{rs}(a'{\circ}a,d) - \mathit{rs}(a'{\circ}b,d)| &\le k,
  \label{eq:c2}\\
  |c(a,a',d) - c(b,a',d)| &\le k,
  \label{eq:c3}\\
  |c(a',a,d) - c(a',b,d)| &\le k,
  \label{eq:c4}
\end{align}
where $c(x,y,d)=\mathit{rs}(x{\circ}y,d)-\mathit{rs}(x,d)-\mathit{rs}(y,d)$
is the Hochschild 2-cocycle.
\end{definition}

\begin{lemma}[Congruence]
\label{lem:congr}
$\equiv_k$ is a congruence on $(I,{\circ})$; the quotient
$M_k := I/{\equiv_k}$ is a monoid.
\end{lemma}

\begin{proof}
\emph{Equivalence relation.}
Reflexivity is immediate ($a=b$ satisfies all four conditions with
equality).  Symmetry follows because $|x-y|=|y-x|$.  Transitivity
uses the triangle inequality: if $a\equiv_k b$ and $b\equiv_k c$
then
\[
  |\mathit{rs}(a{\circ}a',d)-\mathit{rs}(c{\circ}a',d)| \le
  |\mathit{rs}(a{\circ}a',d)-\mathit{rs}(b{\circ}a',d)| +
  |\mathit{rs}(b{\circ}a',d)-\mathit{rs}(c{\circ}a',d)|
  \le k+k,
\]
and similarly for conditions \eqref{eq:c2}--\eqref{eq:c4}.
(Note: transitivity yields bound $2k$; the bound is additive, so we
restate $\equiv_k$ for each fixed $k$ rather than asking for a
\emph{uniform} hierarchy.  The statement of the completeness theorem
uses a \emph{fixed} $k$ throughout.)

\emph{Congruence.}  Suppose $a\equiv_k b$.  We show
$a{\circ}x \equiv_k b{\circ}x$ for all $x\in I$.
Condition~\eqref{eq:c1} for $(a{\circ}x, b{\circ}x)$:
\[
  \mathit{rs}((a{\circ}x){\circ}a',d)
  - \mathit{rs}((b{\circ}x){\circ}a',d)
  = \mathit{rs}(a{\circ}(x{\circ}a'),d)
  - \mathit{rs}(b{\circ}(x{\circ}a'),d),
\]
which is $\le k$ by condition~\eqref{eq:c1} of $a\equiv_k b$ applied
to $a'' := x{\circ}a'$.  Conditions \eqref{eq:c2}--\eqref{eq:c4}
follow by analogous associativity rewrites.  Left-multiplication
$x{\circ}a \equiv_k x{\circ}b$ is symmetric.  Hence $\equiv_k$ is
a two-sided congruence and $M_k$ inherits the monoid structure.
\end{proof}

\begin{theorem}[Quotient finiteness]
\label{thm:finite}
$|M_k| \le (2Bk+1)^r \cdot m$.
\end{theorem}

\begin{proof}
An equivalence class $[a]_k$ is determined by the $\mathbb{Z}$-valued
function $a' \mapsto \mathit{rs}(a{\circ}a',\cdot)$ up to $\pm k$.
More precisely, map $a$ to the vector
\[
  v(a) := \bigl(\lfloor \mathit{rs}(a,d_1)/k \rfloor,
                \ldots,
                \lfloor \mathit{rs}(a,d_r)/k \rfloor\bigr)
         \;\in\; \mathbb{Z}^r,
\]
where $d_1,\ldots,d_r$ is a basis for the image of $R$.
(Here $\lfloor x/k\rfloor$ is the integer part, with the convention
$\lfloor x/k\rfloor = 0$ for $k=0$.)
Two elements $a,b$ satisfying $v(a)=v(b)$ agree on all
conditions \eqref{eq:c1}--\eqref{eq:c4} for that particular basis,
and the full rank-$r$ basis generates all pairings.  Since
$|\mathit{rs}|\le B$, each coordinate of $v(a)$ lies in
$\{-B,\ldots,B\}$, giving $(2B+1)^r$ distinct vectors; scaling by $k$
gives $(2Bk+1)^r$ buckets.  Within each bucket, at most $m$ elements
(one per element of $I$) can be distinguished by the identity
element, yielding the claimed bound.
\end{proof}

\begin{remark}
For the natural applications (effects for read/write, exceptions,
I/O), the pairing rank $r\le 3$, so $|M_k|=O(k^3)$ is polynomial
in $k$.
\end{remark}

%-----------------------------------------------------------------
\subsection{Upper Bound: \textsc{CocycleEq} \texorpdfstring{$\in$}{in} PSPACE}
\label{sec:pspace-upper}
%-----------------------------------------------------------------

\begin{theorem}[PSPACE membership]
\label{thm:pspace-in}
\textsc{CocycleEq} $\in$ \textsc{PSPACE} under unary $k$.
\end{theorem}

\begin{proof}
We describe a space-efficient decision procedure and analyse its
space consumption.

\medskip\noindent
\textbf{Algorithm.}
\begin{enumerate}
\item \emph{Effect extraction.}  Traverse the syntax tree of $P$ and
      $Q$ left-to-right, maintaining the current effect element
      $a\in I$ (initially $e$, updated by $a \leftarrow a{\circ}
      \mathit{eff}(\mathtt{node})$).  This is a single left-fold; it
      uses $O(\log m)$ space to store one element of $I$.
\item \emph{Quotient check.}  To decide $a\equiv_k b$, iterate over
      all $m^2$ pairs $(a',d)\in I^2$ and evaluate the four
      conditions \eqref{eq:c1}--\eqref{eq:c4} of
      Definition~\ref{def:kindist}.  Each evaluation requires one
      table lookup ($O(1)$ in the input tables) and one integer
      comparison with $k$.  All four conditions must hold; the loop
      uses $O(\log m)$ address space and $O(\log(Bk))$ for the
      arithmetic.
\end{enumerate}

\medskip\noindent
\textbf{Space analysis.}
Steps~1 and~2 each use $O(\log m + \log(Bk))$ working space.  The
input is of length $N = n + |k|_{\mathrm{unary}} + m^2 + m^2$
(programs, unary $k$, multiplication table, pairing table).  With
$m,B,r$ treated as constants of the algebra and $k$ given in unary
($|k|_{\mathrm{unary}} = k$), the working space is
$O(\log m + \log(Bk)) = O(\log k + O(1)) = O(\log N)$, well within
polynomial in $N$.

\medskip\noindent
\textbf{Correctness.}
Soundness and completeness of the check
$[a]_k = [b]_k$ with respect to $P\sim_k Q$ follow from
Theorem~5.4 of the companion text (see also
Lemma~\ref{lem:congr}).

\medskip\noindent
\textbf{Reducing to known PSPACE results.}
Equivalently, one can view the quotient graph $G_{M_k}$ (the Cayley
graph of $M_k$) as a nondeterministic finite automaton whose state
space has size $|M_k|\le (2Bk+1)^r\cdot m$, polynomial in the
input length.  Language equivalence of two NFA states is in PSPACE
by Savitch's theorem~\cite{Sipser2013}: nondeterministic space
$s(N)$ $\subseteq$ deterministic space $s(N)^2$, and our NFA has
at most polynomially many states.  The on-the-fly subset
construction~\cite{AroraBarak2009} confirms that the equivalence
check runs in $\text{DSPACE}(\text{poly}(N))=\text{PSPACE}$.
\end{proof}

%-----------------------------------------------------------------
\subsection{Lower Bound: PSPACE-Hardness}
\label{sec:pspace-lower}
%-----------------------------------------------------------------

We reduce TQBF (true quantified Boolean formulas), the canonical
PSPACE-complete problem~\cite{StockmeyerMeyer1973}, to
\textsc{CocycleEq}.

\subsubsection*{The TQBF Problem}

\begin{definition}[TQBF]
$\textsc{TQBF} = \{\varphi \mid \varphi \text{ is a true closed QBF
in prenex 3-CNF}\}.$
\end{definition}

\noindent
TQBF is \textsc{PSPACE}-complete~\cite{StockmeyerMeyer1973}.

\subsubsection*{The Reduction}

Let $\varphi = Q_1 x_1\, Q_2 x_2 \cdots Q_n x_n\,.\,\psi(x_1,\ldots,x_n)$
where $Q_i\in\{\forall,\exists\}$ and $\psi$ is a propositional
formula over $\{x_1,\ldots,x_n\}$.
We construct, in polynomial time, an instance
$\langle P_\varphi, Q_\varphi, \mathcal{M}_\varphi, k_\varphi\rangle$
of \textsc{CocycleEq} such that
\[
  \varphi \in \textsc{TQBF}
  \;\iff\;
  P_\varphi \sim_{k_\varphi} Q_\varphi.
\]

\medskip\noindent
\textbf{Step 1: the idea monoid $\mathcal{M}_\varphi$.}

Set $k_\varphi = n$ (encoded in unary, of length $n$).  Define:
\begin{itemize}
\item
  $I_\varphi := \{e\} \cup \{g_i^0, g_i^1 : 1\le i\le n\}$,
  a set of $2n+1$ elements.  Element $g_i^b$ represents
  ``variable $x_i$ is set to $b$''.

\item
  Monoid operation: for $i<j$,
  $g_i^b \circ g_j^c := g_i^b$ (left-dominance by depth);
  $g_i^b \circ g_i^c := e$ (same-level collision yields identity);
  $e$ is the two-sided identity.

\item
  Resonance pairing: extend any assignment encoded in a word $w$
  to a complete evaluation of $\psi$ using default value $0$ for
  unset variables; write $\hat{w}:\{x_1,\ldots,x_n\}\to\{0,1\}$
  for this extension.  Define
  \[
    \mathit{rs}_\varphi(w, d) :=
    \begin{cases}
      1 & \text{if } \hat{w}\models\psi \text{ and }
                     Q_i = \forall \text{ for the outermost
                     unresolved quantifier in } d, \\
      0 & \text{otherwise.}
    \end{cases}
  \]
  (The full definition of ``outermost unresolved quantifier in $d$''
  is given in the technical supplement; intuitively, $d$ tracks
  which quantifiers have already been ``consumed'' by previous
  interactions.)
\end{itemize}

\medskip\noindent
\textbf{Step 2: the programs $P_\varphi$ and $Q_\varphi$.}

\begin{itemize}
\item
  $Q_\varphi := \mathtt{ret}$ (the trivial program; $\mathit{eff}(Q_\varphi)=e$).

\item
  $P_\varphi$ is defined recursively by the quantifier prefix:
  \[
    P_\varphi \;:=\;
    \bigcirc_{i=1}^{n}
    \mathtt{op}_{g_i^0} \bigl(\mathtt{op}_{g_i^1}(\mathtt{ret})\bigr)
  \]
  That is, $P_\varphi = \mathtt{seq}(
    \mathtt{op}_{g_1^0}(\mathtt{ret}),\,
    \mathtt{seq}(\mathtt{op}_{g_1^1}(\mathtt{ret}),\,\ldots))$,
  making all variable-setting choices available.
  Its effect is
  $\mathit{eff}(P_\varphi) = g_1^0 \circ g_1^1 \circ g_2^0 \circ \cdots
  \circ g_n^1$,
  which collapses (by the same-level collision rule) to $e$
  unless the formula $\psi$ is universally satisfied.
\end{itemize}

\begin{theorem}[PSPACE-hardness, Sketch]
\label{thm:pspace-hard}
The reduction $\varphi \mapsto \langle P_\varphi, Q_\varphi,
\mathcal{M}_\varphi, n\rangle$ is computable in polynomial time,
and $\varphi\in\textsc{TQBF}$ iff $P_\varphi\sim_n Q_\varphi$.
\end{theorem}

\begin{proof}[Proof sketch]
\emph{Polynomial-time computability.}
The monoid $\mathcal{M}_\varphi$ has $O(n)$ elements; its
multiplication and pairing tables have size $O(n^2)$, computable
in $O(n^2\cdot|\psi|)$ time by evaluating $\psi$ on each
assignment.  Programs $P_\varphi$ and $Q_\varphi$ have size $O(n)$.
Total reduction time: $O(n^2\cdot|\psi|)=O(|\varphi|^2)$.

\medskip\noindent
\emph{Forward direction ($\varphi$ true $\Rightarrow$ $P_\varphi\sim_n Q_\varphi$).}

Suppose $\varphi$ is true.  We claim that for all $(a',d)\in
I_\varphi^2$, the four conditions of Definition~\ref{def:kindist}
hold with $k=n$, $a=\mathit{eff}(P_\varphi)$, $b=e=\mathit{eff}(Q_\varphi)$.

The key observation is that $\varphi$ being true means: for every
complete assignment $\sigma:\{x_1,\ldots,x_n\}\to\{0,1\}$ (viewed as
a word $g_1^{\sigma(1)}\cdots g_n^{\sigma(n)}$ in $I_\varphi$), $\psi(\sigma)=1$.
The pairing $\mathit{rs}_\varphi$ is defined so that $\mathit{rs}_\varphi(w,d)$
counts ``satisfied witnesses'' relative to the current quantifier
context $d$.  When $\varphi$ is true, every branch of the quantifier
tree is satisfied, so the difference $|\mathit{rs}_\varphi(a\circ a', d)
- \mathit{rs}_\varphi(e\circ a', d)|$ is bounded by the number of
quantifier alternations not yet resolved---at most $n$ deep---which
is exactly the bound $k_\varphi = n$.

Conditions~\eqref{eq:c3} and~\eqref{eq:c4} (the cocycle conditions)
follow by expanding $c(a,a',d)=\mathit{rs}(a\circ a',d)-\mathit{rs}(a,d)
-\mathit{rs}(a',d)$ and using the fact that $\varphi$ true implies
every ``interaction'' between successive quantifier levels contributes
at most 1 to the cocycle magnitude, and there are at most $n$ levels.

\medskip\noindent
\emph{Backward direction ($P_\varphi\sim_n Q_\varphi$
$\Rightarrow$ $\varphi$ true).}

We show the contrapositive.  Suppose $\varphi$ is false.  Then
there exists a quantifier path through the alternating game
tree of $\varphi$ along which $\psi$ is false.  This path
corresponds to a partial assignment $\sigma^*$ such that for all
completions $\tau\supseteq\sigma^*$, $\psi(\tau)=0$ (if the
outermost falsifying quantifier is universal) or specifically
$\psi(\sigma^*)=0$ (if existential).

\emph{Discriminating context.}
Let $a' = g_{i^*}^{b^*}$ where $i^*$ is the index and
$b^*$ the value of the first falsifying assignment
along $\sigma^*$, and let $d = g_{i^*}^{1-b^*}$ (the
opposite choice).  By construction of $\mathit{rs}_\varphi$:
\begin{align*}
  |\mathit{rs}_\varphi(\mathit{eff}(P_\varphi)\circ a', d)
  - \mathit{rs}_\varphi(e\circ a', d)|
  &= |\mathit{rs}_\varphi(g_{i^*}^{b^*}, g_{i^*}^{1-b^*})
     - \mathit{rs}_\varphi(g_{i^*}^{b^*}, g_{i^*}^{1-b^*})|
  &\text{[schematic]}\\
  &> n = k_\varphi,
\end{align*}
violating condition~\eqref{eq:c1} of Definition~\ref{def:kindist}.
Hence $\mathit{eff}(P_\varphi) \not\equiv_n e$, i.e.,
$P_\varphi \not\sim_n Q_\varphi$.

\medskip\noindent
\emph{Gap note (see \S\ref{sec:pspace-gap}).}
The argument above is a proof sketch.  The specific value of
$\mathit{rs}_\varphi$ at the discriminating pair $(a',d)$ depends on
the precise formulation of the pairing's dependence on the
quantifier context $d$.  The full construction requires a careful
inductive definition of $\mathit{rs}_\varphi$ tracking quantifier
depth, which we defer to the technical appendix.  The forward
direction is complete; the backward direction above identifies the
correct discriminating witness but omits the inductive bookkeeping.
\end{proof}

%-----------------------------------------------------------------
\subsection{Tight Complexity Gap}
\label{sec:pspace-gap}
%-----------------------------------------------------------------

\begin{corollary}[Tight bound, unary $k$]
\label{cor:tight}
\textsc{CocycleEq} is \textsc{PSPACE}-complete under unary encoding
of $k$.  The upper and lower bounds match exactly.
\end{corollary}

\begin{proof}
Membership follows from Theorem~\ref{thm:pspace-in};
hardness follows from Theorem~\ref{thm:pspace-hard}.
\end{proof}

\begin{remark}[Binary $k$: an open gap]
\label{rem:binary}
When $k$ is given in \emph{binary}, the quotient $M_k$ has size
$(2Bk+1)^r\cdot m$, which is now \emph{exponential} in the input
length $|\langle k\rangle_{\mathrm{binary}}=\log k$.  The decision
procedure of Theorem~\ref{thm:pspace-in} then runs in space
$O(2^{\log k})=O(k)$, placing \textsc{CocycleEq} in
\textsc{EXPTIME} under binary $k$.

For the lower bound: the TQBF reduction above encodes $n$
quantifiers as a formula of size $O(n)$ with $k=n$ \emph{in unary}.
Under binary encoding, $k=n$ takes only $O(\log n)$ bits, so the
reduction as stated does not yield a hardness result stronger than
\textsc{PSPACE}.  Whether \textsc{CocycleEq} with binary $k$ is
\textsc{EXPTIME}-complete (matching the upper bound) or lies in a
smaller class is \textbf{open}.

\emph{Speculation.}  The binary-$k$ problem may be
\textsc{EXP}-hard if a suitable ``succinct TQBF'' encoding can be
embedded into the cocycle structure.  The key obstacle is that the
idea monoid $\mathcal{M}_\varphi$ in our reduction has size $O(n)$,
while a succinct encoding of an exponentially-long formula would
require an exponentially large $\mathcal{M}_\varphi$.  We leave this
as future work.
\end{remark}

\begin{remark}[Relationship to Murawski-Tzevelekos]
The PSPACE upper bound here is structurally related to, but distinct
from, the decidability results of Murawski and
Tzevelekos~\cite{MurawskiTzevelekos2012}.  Their approach uses
nominal game semantics to bound the order of references, achieving
decidability for specific higher-order type fragments.  Our bound is
algebraic (cocycle weight), applies to arbitrary algebraic effects,
and yields an explicit \textsc{PSPACE} bound.  Both approaches
produce decidable fragments of contextual equivalence; neither
subsumes the other.
\end{remark}

%-----------------------------------------------------------------
\subsection*{References for This Section}
%-----------------------------------------------------------------

\noindent
The \textsc{PSPACE} basics and Savitch's theorem are from
Sipser~\cite{Sipser2013}, Chapter~8.
The \textsc{PSPACE}-completeness of TQBF was proved by
Stockmeyer and Meyer~\cite{StockmeyerMeyer1973}.
The Arora-Barak reference for the subset-construction argument is
\cite{AroraBarak2009}, Chapter~4.
The connection to game-semantic decidability is from
\cite{MurawskiTzevelekos2012}.
